\SourceLeaf{222}%
\bigskip\Needspace{10\baselineskip}
\pdfbookmark[1]{Further Report. Pallas Olbersiana. July 1802}{pallas180207}
\begin{center}
{\small FURTHER REPORT ON THE NEW PRIMARY PLANET\par}
PALLAS OLBERSIANA.\par
\medskip\hrule\medskip
Monatliche Correspondenz. Vol. VI. pp.\ 71—96. July 1802.\par
\medskip\hrule
\end{center}\medskip
—\quad —\quad —

“The very instructive paper by Prof.\ \textsc{Wurm} in your June issue (Dr.\ \textsc{Gauss} writes on 26 June) unfortunately shows that, given the uncertainty that still surrounds Jupiter’s mass, we will be unable as yet to determine the perturbation equations of \textit{Ceres} and \textit{Pallas} with the reliability one might wish for. On the other hand, the observations of \textit{Ceres}, and still more those of \textit{Pallas}, will also be of great use because they will allow that mass to be determined very well after one or a couple of revolutions. That in turn can only have a beneficial effect on many other important points in astronomy. Thus in this sublime science, which heaven seems to have bestowed on humankind expressly to cultivate its powers and lift it above earthly things, one discovery always holds out its hand to another. Great discoveries never stand alone and are often as important in their consequences as they are in themselves. That very uncertainty will give me another reason to keep to the intention I mentioned earlier. Besides the new ellipse for \textit{Ceres}, which is at least to be calculated with the closest attention to Jupiter’s perturbations, I intend to calculate another, \textit{pure} ellipse, based on exactly the same observations and free of all perturbations. I think it would be interesting to see to what extent a \textit{definite} trace of Jupiter’s effects will already appear in the motion of \textit{Ceres} next year. For the present there is not the slightest trace, and the observations of a year and a half will no doubt be represented quite well by a \textit{pure ellipse} without forcing the fit. That elements VII are off by $40''$ in May need not surprise anyone at all, considering how scanty the underlying observations were. After all, I had only the Palermo observations and a couple from you for which the declinations were missing. For that reason I still cannot abandon my earlier opinion that taking the perturbations into account even then would have added quite needlessly to the work. I really do not see how it could have made anything better up to now, but I do believe that it would have made a great deal worse. I at least would hardly have had the patience to compare so many observations with the elements down to a hair and to calculate an ephemeris for a considerable time, if I had also had to calculate so many equations every time. An ellipse calculated without regard to the perturbations, which fits the observations exactly, must itself include those perturbations for a considerable time. As long as our only concern is to facilitate observations, examine those already made and assess their respective worth, this ellipse must be far more important to us than a \textit{mean ellipse}. That will only become interesting once the observations span enough time for it to be determined with sufficient certainty to go back a number of years with it and search the older star catalogues. After a year and a half one can now already expect something fairly reliable from it. Accordingly, once I have the last of the Palermo
\SourceLeaf{223}%
observations, I will carry out these calculations with all the care I can. For these reasons I also think that there would be little benefit in taking perturbations into account for \textit{Pallas} as early as this year, especially because in my opinion the methods used hitherto for other planets would by no means give the required precision for this one. Only if the faintness of the planet, \textit{quod Deus avortat}, should prevent its being found again next year, so that we would have to project two years ahead with this year’s observations, would it be advisable to take Jupiter’s effect into account, though, in my view, by a method other than the usual one.”\quad —\quad —\quad —

\Needspace{14\baselineskip}
\begin{center}
\textit{Elements of Pallas Olbersiana from 42 days of meridian observations at Seeberg and Greenwich.}\par\medskip
(III.)\par\medskip
{\small\setlength{\tabcolsep}{2pt}
\begin{tabular}{l r}
Epoch 1802 March 31, noon at Seeberg $\ .\ .\ .\ .\ .\ .\ .$&$162^0\ 25'\ 45''9$\\
mean daily tropical motion $\ .\ .\ .\ .\ .\ .\ .$&$769''547$\\
Logarithm of the semimajor axis $\ .\ .\ .\ .\ .\ .\ .\ .$&$0.4425664$
\end{tabular}
\par\smallskip
$\left.\begin{array}{l}\text{Aphelion}\\\text{Node}\quad .\end{array}\right\}\ \text{for the epoch and fixed relative to the stars}\left\{\begin{array}{lr}
.\ .&300^0\ 58'\ 47''7\\
.\ .&172^0\ 28'\ 17''9
\end{array}\right.$
\par\smallskip
\begin{tabular}{l r}
Eccentricity $\ .\ .\ .\ .\ .\ .\ .\ .\ .\ .\ .\ .\ .\ .\ .\ .\ .$&$0.2476402$\\
Inclination $\ .\ .\ .\ .\ .\ .\ .\ .\ .\ .\ .\ .\ .\ .\ .\ .\ .\ .$&$34^0\ 39'\ 10''7$
\end{tabular}}
\end{center}

It is indeed most remarkable that the mean tropical motion of \textit{Pallas} is coming ever closer to that of \textit{Ceres}, without anyone intending it and entirely of its own accord. From elements IV of the orbit of Ceres (December 1801 issue, p.\ 647 [203 of this volume]), Dr.\ \textsc{Gauss} determined this mean motion to be $770''914$. It differs from that given above for \textit{Pallas} by only one and a half seconds. But since some uncertainty still attaches to both motions, one cannot really state any difference at all as yet.

\Needspace{27\baselineskip}
Dr.\ \textsc{Gauss} then compared the following series of observations with elements III above and found the splendid agreement shown in the following table.
\begin{center}
{\small\setlength{\tabcolsep}{4pt}
\begin{tabular}{r r|r|r|r|r|l}
\multicolumn{2}{c|}{1802}&\shortstack{Calculated\\right ascension\\of Pallas}&\shortstack{Calculated\\declination\\of Pallas}&\multicolumn{2}{c|}{Difference}&Observer.\\
\multicolumn{2}{c|}{}&&&in R.A.&in decl.&\\\hline
April&4&$183^0\ 44'\quad5''6$&$13^0\ 54'\ 53''5$&$-1''0$&$+1''5$&von Zach\\
&5&$183\ 34\ 24.4$&$14\ 13\ 20.2$&$+0.$&$-2.7$&von Zach\\
&7&$183\ 15\ 38.5$&$14\ 48\ 59.3$&$0.7$&$-2.8$&von Zach\\
&8&$183\quad6\ 33.9$&$15\quad6\ 12.8$&$-3.9$&$+2.8$&von Zach\\
&10&$182\ 48\ 56.3$&$15\ 39\ 43.5$&$-1.7$&$+1.5$&Dr.\ Burckhardt\\
&12&$182\ 32\ 33.7$&$16\ 11\quad6.4$&$+2.2$&$+8.4$&Méchain\\
&13&$182\ 24\ 46.8$&$16\ 26\quad8.3$&$+2.4$&$+2.3$&Méchain\\
&15&$182\ 10\ 14.1$&$16\ 54\ 33.7$&$-2.4$&$+2.9$&von Zach\\
&15&$182\ 10\quad4.2$&$16\ 54\ 53.3$&$+1.3$&$+1.9$&Méchain\\
&16&$182\quad3\quad9.7$&$17\quad8\ 36.3$&$+1.7$&$-0.5$&Méchain\\
&17&$181\ 56\ 33.7$&$17\ 21\ 52.8$&$+3.4$&$+3.0$&Méchain\\
&18&$181\ 50\ 25.0$&$17\ 34\ 25.3$&$-5.6$&$\ .\ .\ .$&von Zach\\
&19&$181\ 44\ 26.4$&$17\ 46\ 50.0$&$+1.1$&$-4.4$&von Zach\\
&23&$181\ 23\ 41.9$&$18\ 32\ 27.4$&$-0.1$&$-0.6$&Dr.\ Maskelyne\\
&24&$181\ 19\ 32.4$&$18\ 42\ 25.2$&$-5.4$&$\ .\ .\ .$&von Zach\\
&25&$181\ 15\ 35.3$&$18\ 52\ 15.9$&$+3.1$&$\ .\ .\ .$&von Zach\\
&25&$181\ 15\ 28.6$&$18\ 52\ 33.1$&$+0.3$&$+3.4$&Dr.\ Maskelyne\\
&26&$181\ 11\ 59.4$&$19\quad1\ 41.6$&$+2.4$&$-7.8$&von Zach\\
&26&$181\ 11\ 53.3$&$19\quad1\ 58.1$&$+3.1$&$-4.5$&Dr.\ Maskelyne\\
&27&$181\quad8\ 44.8$&$19\ 10\ 42.7$&$+0.8$&$-4.1$&von Zach
\end{tabular}}
\end{center}
\expandafter\def\csname HistoricalNoteText20\endcsname{There seems to have been some oversight in the \textit{Greenwich} observation of the R.A. on 7 May, since it seems to be about $39''$ too small.}
\SourceLeaf{224}%
\begin{center}
{\small\setlength{\tabcolsep}{4pt}
\begin{tabular}{r r|r|r|r|r|l}
\multicolumn{2}{c|}{1802}&\shortstack{Calculated\\right ascension\\of Pallas}&\shortstack{Calculated\\declination\\of Pallas}&\multicolumn{2}{c|}{Difference}&Observer\\
\multicolumn{2}{c|}{}&&&in R.A.&in decl.&\\\hline
April&29&$181^0\quad3'\ 20''9$&$19^0\ 27'\ 31''5$&$+4''3$&$-12''2$&von Zach\\
&30&$181\quad1\ 12.1$&$19\ 35\ 19.6$&$+2.2$&$-13.9$&von Zach\\
May&1&$180\ 59\ 24.6$&$19\ 42\ 44.8$&$+2.1$&$+4.0$&von Zach\\
&1&$180\ 59\ 21.7$&$19\ 42\ 57.7$&$+2.8$&$+3.3$&Dr.\ Maskelyne\\
&2&$180\ 57\ 59.7$&$19\ 49\ 46.5$&$+7.8$&$-5.5$&von Zach\\
&2&$180\ 57\ 57.5$&$19\ 49\ 58.7$&$+0.2$&$+2.0$&Dr.\ Maskelyne\\
&3&$180\ 56\ 56.9$&$19\ 56\ 25.9$&$+1.3$&$-21.8$&von Zach\\
&4&$180\ 56\ 15.0$&$20\quad2\ 53.0$&$+2.1$&$-0.3$&Dr.\ Maskelyne\\
&5&$180\ 55\ 56.9$&$20\quad8\ 36.2$&$+2.3$&$-2.5$&von Zach\\
&6&$180\ 56\quad0.3$&$20\ 14\quad8.8$&$+6.2$&$\ .\ .\ .$&von Zach\\
&7&$180\ 56\ 24.7$&$20\ 19\ 20.4$&$+0.9$&$-14.9$&von Zach\\
&7&$180\ 56\ 25.8$&$20\ 19\ 29.4$&$+39.3$&$-0.5$&Dr.\ Maskelyne\\
&8&$180\ 57\ 13.0$&$20\ 24\ 10.8$&$+3.0$&$-9{,}3$&von Zach\\
&9&$180\ 58\ 23.9$&$20\ 28\ 49.3$&$+2.3$&$+6.8$&Dr.\ Maskelyne\\
&11&$181\quad1\ 44.6$&$20\ 36\ 41.9$&$+2.9$&$-31.3$&von Zach\\
&11&$181\quad1\ 48.3$&$20\ 36\ 48.5$&$+3.4$&$+0.3$&Dr.\ Maskelyne\\
&13&$181\quad6\ 37.6$&$20\ 43\ 31.4$&$+4.6$&$-1.7$&Dr.\ Maskelyne\\
&14&$181\quad9\ 33.4$&$20\ 46\ 25.2$&$+4.0$&$-1.5$&Dr.\ Maskelyne\\
&16&$181\ 16\ 26.5$&$20\ 51\ 19.3$&$-1.7$&$+0.5$&Dr.\ Maskelyne
\end{tabular}}
\end{center}

This comparison shows that, on the basis of this entire series of observations, little further change could be made to the elements, especially since the last Greenwich declinations in May (which have a greater influence on the elements than the right ascensions) still agree very well with one another. Only the calculated right ascensions\HistoricalNote{20} would need to be reduced by about $4''$ toward the end, and that would have only a negligible effect on the elements. In a letter of 26 June, Dr.\ \textsc{Gauss} informs us that he had also compared \textsc{Olbers}’s observations of this planet on 19, 20 and 21 June with his elements III. They still agreed so beautifully that he could make no change at all with confidence on their basis. We can therefore keep these orbital elements III until the observations are fully concluded. Dr.\ \textsc{Gauss} accordingly used them to calculate the following
\Needspace{15\baselineskip}
\begin{center}
\textit{Ephemeris for Pallas Olbersiana for midnight at Seeberg.}\par\medskip
{\small\setlength{\tabcolsep}{3pt}
\begin{tabular}{r r|r|r|r|r}
\multicolumn{2}{c|}{1802}&\shortstack{R.A.\\of Pallas}&\shortstack{Decl.\\of Pallas}&\shortstack{Distance\\from the $\EarthSymbol$}&\shortstack{Distance\\from the $\odot$}\\\hline
May&18&$181^0\ 25'$&$20^0\ 55'$&$1.888$&$2.477$\\
&21&$181\ 40$&$20\ 59$&$1.930$&$2.484$\\
&24&$181\ 58$&$21\quad0$&$1.973$&$2.492$\\
&27&$182\ 19$&$20\ 59$&$2.016$&$2.500$\\
&30&$182\ 42$&$20\ 56$&$2.060$&$2.508$\\
June&2&$183\quad8$&$20\ 51$&$2.104$&$2.516$\\
&5&$183\ 36$&$20\ 44$&$2.148$&$2.524$\\
&8&$184\quad6\frac{1}{2}$&$20\ 36$&$2.192$&$2.532$
\end{tabular}
\par\medskip
\begin{tabular}{r r|r|r|r|r}
\multicolumn{2}{c|}{1802}&\shortstack{R.A.\\of Pallas}&\shortstack{Decl.\\of Pallas}&\shortstack{Distance\\from the $\EarthSymbol$}&\shortstack{Distance\\from the $\odot$}\\\hline
June&11&$184^0\ 39'$&$20^0\ 26\frac{1}{2}'$&$2.237$&$2{,}539$\\
&14&$185\ 14$&$20\ 15$&$2.281$&$2.547$\\
&17&$185\ 51$&$20\quad3$&$2.326$&$2.555$\\
&20&$186\ 30$&$19\ 50$&$2.370$&$2.563$\\
&23&$187\ 10$&$19\ 35$&$2.415$&$2.571$\\
&26&$187\ 52\frac{1}{2}$&$19\ 19\frac{1}{2}$&$2.459$&$2.579$\\
&29&$188\ 36$&$19\quad3$&$2.504$&$2.587$
\end{tabular}}
\end{center}

For the beginning of the coming year, Dr.\ \textsc{Gauss} calculated the approximate position of \textit{Pallas} as follows. On 31 Dec.\ at $12^{\mathrm{u}}$ midnight at Seeberg, longitude $253^0\ 33'$, latitude $26^0\ 10'$, N., distance from the Earth $3.779$, from the Sun $3.043$. Given the faintness of this planet, however, it remains very uncertain whether we will see it again
\SourceLeaf{225}%
in the next two years. Dr.\ \textsc{Olbers} calculated the next two oppositions of \textit{Pallas} from \textsc{Gauss}’s elements II. When he discovered it on 28 March, its distance from the Earth was only $=1.38$. At the end of May it was $=2.00$, and on 15 June $=2.31$. Next year, in 1803, toward the end of June, \textit{Pallas} will again be in opposition to the Sun. It will be above Cerberus at longitude $9^{\mathrm{z}}\ 8^0$ and $48^0$ north latitude, but its distance from the Earth will be $2.61$. Dr.\ \textsc{Olbers} thinks that it will hardly be possible to find it at this distance in the bright northern twilight. The opposition of 1804 is somewhat more favorable. It occurs toward the end of August at $11^{\mathrm{z}}\ 3^0$ with latitude $19^0$ N., thus in the head of Pegasus. Dr.\ \textsc{Olbers} still finds the distance from the Earth to be $2.40$, but the nights are already completely dark. Dr.\ \textsc{Olbers} believes, however, that we will \textit{most certainly} see \textit{Pallas} again in 1805, since at opposition it will already have descended into the part of its orbit nearer the Sun and thus will also be much nearer the Earth. Meanwhile, we hope these considerations will deter no astronomer from doing their utmost to track down \textit{Pallas} as soon as a search is possible. For large telescopes with high light-gathering power might still be enough to detect this planet, even if it appears only as a star of the twelfth magnitude. Finding it will of course be somewhat more difficult. But even supposing that all attempts fail in 1803 and 1804, Dr.\ \textsc{Gauss} still hopes that his orbital elements for 1805 will be a sufficient guide for finding \textit{Pallas} again (with somewhat more trouble, admittedly, than in the case of \textit{Ceres}).\quad —\quad —\quad —

\Needspace{30\baselineskip}
As for the small size of these two new planetary bodies, it is indeed difficult to see how so small an apparent diameter can be measured accurately. \textsc{Herschel} and \textsc{Schröter}, both practiced observers equipped with the largest and best optical instruments, have attempted such measurements. Schröter found \textit{Ceres} 18 times and \textit{Pallas} 40 times as large as Herschel did. To show the difficulty of such measurements, Dr.\ \textsc{Gauss} used his latest orbital elements to make the following calculations of the apparent diameters of these bodies.
\begin{center}
\textit{Diameter of Ceres.}\par\medskip
{\small\setlength{\tabcolsep}{6pt}
\begin{tabular}{r r|r|c}
\multicolumn{2}{c|}{1802}&\shortstack{Distance\\from the\\$\EarthSymbol$}&\shortstack{Apparent diameter, taking the true diameter,\\according to Dr.\ \textsc{Herschel}, as $=35$ German miles.}\\\hline
March&28&$1.621$&$0''22$
\end{tabular}}
\par\bigskip
\textit{Diameter of Pallas.}\par\medskip
{\small\setlength{\tabcolsep}{6pt}
\begin{tabular}{r r|r|c}
\multicolumn{2}{c|}{1802}&\shortstack{Distance\\from the\\$\EarthSymbol$}&\shortstack{Apparent diameter, assuming that\\Dr.\ \textsc{Herschel}'s result $=15$ German\\miles is correct.}\\\hline
April&4&$1.407$&$0''11$\\
—&15&$1.386$&$0.10$\\
June&23&$2.415$&$0.06$
\end{tabular}}
\end{center}

We leave it to others to judge whether such small diameters can actually be measured, or whether one can even see the light reflected from such small surfaces.\quad —\quad —\quad $=$
\begin{flushright}[\textsc{von Zach}.]\end{flushright}
